\documentclass[11pt]{article}
\usepackage[utf8]{inputenc}
\usepackage{amsmath,amssymb}
\usepackage{hyperref}
\title{Efficient Single Atom Catalyst: A Resource-Aware Approach}
\author{Anonymous}
\date{}
\begin{document}
\maketitle

\begin{abstract}
This paper studies Single Atom Catalyst. Grounded in a real, citable literature \cite{ref1,ref2,ref3,ref4,ref5,ref6,ref7,ref8},
we propose a focused method and report \textbf{illustrative} experiments.
All references below are real and verifiable (OpenAlex/arXiv); the reported
numbers are simulated to illustrate intended behaviour.
\end{abstract}

\section{Introduction}
Single Atom Catalyst is an active research area. This work targets a concrete gap identified from
the recent literature and contributes a method together with an evaluation protocol.

\section{Related Work (real literature)}
The following works, retrieved from public bibliographic APIs, frame our study:
\begin{itemize}
\item Qiao et al. (2011) — "Single-atom catalysis of CO oxidation using Pt1/FeOx", Nature Chemistry (cited 7373).
\item Yang et al. (2013) — "Single-Atom Catalysts: A New Frontier in Heterogeneous Catalysis", Accounts of Chemical Research (cited 4719).
\item Liu et al. (2018) — "Metal Catalysts for Heterogeneous Catalysis: From Single Atoms to Nanoclusters and Nanoparticles", Chemical Reviews (cited 4706).
\item Wang et al. (2018) — "Heterogeneous single-atom catalysis", Nature Reviews Chemistry (cited 4195).
\item Batzill et al. (2005) — "The surface and materials science of tin oxide", Progress in Surface Science (cited 2521).
\item Yin et al. (2016) — "Single Cobalt Atoms with Precise N‐Coordination as Superior Oxygen Reduction Reaction Catalysts", Angewandte Chemie International Edition (cited 2251).
\end{itemize}

\section{Method}
We reduce the dominant computational cost via adaptive resource allocation, activating only the components needed per input. We instantiate this idea for Single Atom Catalyst and analyse its cost/quality trade-off.

\section{Experiments (Illustrative)}
\begin{center}
\begin{tabular}{lcc}
System & Primary Metric & Efficiency \\ \hline
Strong baseline & 0.812 & 1.00$\times$ \\
Ablation & 0.828 & 0.92$\times$ \\
\textbf{Ours} & \textbf{0.851} & \textbf{0.71$\times$}
\end{tabular}
\end{center}
\emph{Metrics simulated to illustrate the intended effect; not measured.}

\section{Conclusion}
We connected a real literature to a concrete method for Single Atom Catalyst. Future work will
validate the illustrative results with real experiments.

\begin{thebibliography}{9}
\bibitem{ref1} Botao Qiao, Aiqin Wang, Xiaofeng Yang et al.. Single-atom catalysis of CO oxidation using Pt1/FeOx. \emph{Nature Chemistry}, 2011. DOI:10.1038/nchem.1095
\bibitem{ref2} Xiao-Feng Yang, Aiqin Wang, Botao Qiao et al.. Single-Atom Catalysts: A New Frontier in Heterogeneous Catalysis. \emph{Accounts of Chemical Research}, 2013. DOI:10.1021/ar300361m
\bibitem{ref3} Lichen Liu, Avelino Corma. Metal Catalysts for Heterogeneous Catalysis: From Single Atoms to Nanoclusters and Nanoparticles. \emph{Chemical Reviews}, 2018. DOI:10.1021/acs.chemrev.7b00776
\bibitem{ref4} Aiqin Wang, Jun Li, Tao Zhang. Heterogeneous single-atom catalysis. \emph{Nature Reviews Chemistry}, 2018. DOI:10.1038/s41570-018-0010-1
\bibitem{ref5} Matthias Batzill, Ulrike Diebold. The surface and materials science of tin oxide. \emph{Progress in Surface Science}, 2005. DOI:10.1016/j.progsurf.2005.09.002
\bibitem{ref6} Peiqun Yin, Tao Yao, Yuen Wu et al.. Single Cobalt Atoms with Precise N‐Coordination as Superior Oxygen Reduction Reaction Catalysts. \emph{Angewandte Chemie International Edition}, 2016. DOI:10.1002/anie.201604802
\bibitem{ref7} Yuanjun Chen, Shufang Ji, Chen Chen et al.. Single-Atom Catalysts: Synthetic Strategies and Electrochemical Applications. \emph{Joule}, 2018. DOI:10.1016/j.joule.2018.06.019
\bibitem{ref8} Hong Bin Yang, Sung‐Fu Hung, Song Liu et al.. Atomically dispersed Ni(i) as the active site for electrochemical CO2 reduction. \emph{Nature Energy}, 2018. DOI:10.1038/s41560-017-0078-8
\end{thebibliography}
\end{document}
